% ============================================================
%  CHAPTER 5 : CHALLENGES AND THE BLOCKCHAIN TRILEMMA
% ============================================================
\chapter{Challenges in Blockchain and the Blockchain Trilemma}
\label{ch:challenges}

\noindent\textbf{Declaration of Contributions:} This chapter was primarily authored by \textbf{Madhav Deorah}, who researched and wrote the sections on scalability issues, energy consumption, latency, transaction fees, and the Blockchain Trilemma. \textbf{Anmol Goklani} contributed to the formalisation of the trilemma trade-off. All team members reviewed the chapter.

\bigskip

Despite its transformative potential, blockchain technology faces several fundamental challenges that must be addressed before it can achieve mainstream adoption. This chapter examines these challenges and formalises the central trade-off that underlies many of them: the Blockchain Trilemma.

\section{Scalability Issues}
\label{sec:scalability-challenge}

As discussed in \Cref{sec:scalability}, blockchain networks process far fewer transactions per second than traditional systems. Bitcoin handles approximately 7 \gls{tps}, Ethereum approximately 15--30 \gls{tps}, while Visa processes around 50{,}000 \gls{tps}~\cite{kohad2020scalability}.

This limitation becomes increasingly problematic as blockchain usage grows. During periods of high demand, networks become congested, transaction confirmation times increase, and fees rise sharply. For example, during the 2017 and 2021 cryptocurrency booms, Ethereum transaction fees sometimes exceeded \$50--\$100 per transaction.

The root cause is architectural: in most blockchain designs, \emph{every node processes every transaction}. This provides maximum security and decentralisation but creates an inherent throughput ceiling.

\section{High Energy Consumption}
\label{sec:energy}

\Gls{pow} consensus, as used by Bitcoin, requires miners to perform vast amounts of computation. The Bitcoin network's total energy consumption is estimated at 120--180 TWh per year---comparable to the annual electricity usage of countries like Argentina or Norway.

This has drawn significant criticism from environmentalists and policymakers. The transition of Ethereum from PoW to \gls{pos} in September 2022 (``The Merge'') demonstrated that energy-efficient alternatives exist, reducing Ethereum's energy consumption by approximately 99.95\%.

\section{Latency and Speed}
\label{sec:latency}

Blockchain transactions are not instant. Before a transaction is considered ``confirmed,'' it must be:
\begin{enumerate}
    \item Broadcast to the network.
    \item Included in a block by a miner/validator.
    \item Followed by several additional blocks (to reduce the risk of reorganisation).
\end{enumerate}

In Bitcoin, a single block takes approximately 10 minutes to mine, and 6 confirmations (approximately 1 hour) are typically recommended for high-value transactions. Even Ethereum, with its faster 12-second block time, requires multiple confirmations for finality.

This latency makes blockchain unsuitable for applications requiring real-time transaction processing, such as point-of-sale payments.

\section{High Transaction Fees}
\label{sec:fees}

Transaction fees in blockchain networks are market-driven: when demand for block space exceeds supply, users must bid higher fees to have their transactions included promptly. During network congestion, this creates a ``fee market'' where:
\begin{itemize}
    \item Wealthy users can afford fast confirmation.
    \item Ordinary users face long delays or prohibitive costs.
    \item Small-value transactions become economically unviable.
\end{itemize}

This is particularly problematic for the vision of blockchain as a tool for \emph{financial inclusion}---providing banking services to the unbanked.

\section{The Blockchain Trilemma}
\label{sec:trilemma}

The challenges above are not independent; they are manifestations of a fundamental trade-off first articulated by Ethereum co-founder Vitalik Buterin. The \textbf{Blockchain Trilemma} states that a blockchain system can optimise at most two of the following three properties simultaneously~\cite{dai2019trilemma}:

\begin{enumerate}
    \item \textbf{Decentralisation:} The network is controlled by many independent nodes, with no single point of authority.
    \item \textbf{Security:} The network is resistant to attacks and manipulation.
    \item \textbf{Scalability:} The network can process a large volume of transactions quickly and cheaply.
\end{enumerate}

\begin{figure}[htbp]
\centering
\begin{tikzpicture}[
    vertex/.style={circle, draw, thick, minimum size=2.2cm, font=\small\bfseries, align=center},
    >=Stealth
]

\node[vertex, fill=blue!12] (D) at (90:3.5) {Decentra-\\lisation};
\node[vertex, fill=red!12] (Se) at (210:3.5) {Security};
\node[vertex, fill=green!12] (Sc) at (330:3.5) {Scalability};

\draw[very thick, gray!50] (D) -- (Se);
\draw[very thick, gray!50] (Se) -- (Sc);
\draw[very thick, gray!50] (Sc) -- (D);

\node[font=\large\bfseries, text=gray!70] at (0, 0) {Pick Two};

\node[below=0.3cm of D, font=\scriptsize, text width=3.5cm, align=center] {Distributed control\\No central authority};
\node[below left=0.3cm of Se, font=\scriptsize, text width=3cm, align=center] {Attack resistance\\Network trust};
\node[below right=0.3cm of Sc, font=\scriptsize, text width=3.5cm, align=center] {High throughput\\Fast processing};

\end{tikzpicture}
\caption{The Blockchain Trilemma. Any blockchain system can optimise at most two of the three properties---decentralisation, security, and scalability---simultaneously. Improving one typically comes at the expense of another.}
\label{fig:trilemma}
\end{figure}

\Cref{tab:trilemma-examples} illustrates how existing systems navigate this trade-off.

\begin{table}[htbp]
\centering
\renewcommand{\arraystretch}{1.3}
\begin{tabularx}{\textwidth}{l c c c X}
\toprule
\textbf{System} & \textbf{Decentralised} & \textbf{Secure} & \textbf{Scalable} & \textbf{Trade-off} \\
\midrule
Bitcoin       & \checkmark & \checkmark & $\times$   & Low throughput (7 TPS) \\
Solana        & $\times$   & \checkmark & \checkmark & Fewer validators (centralisation) \\
Hyperledger   & $\times$   & \checkmark & \checkmark & Permissioned (not open) \\
Ethereum 2.0  & \checkmark & \checkmark & Improving  & Still working on full scaling \\
\bottomrule
\end{tabularx}
\caption{How different blockchain systems navigate the Blockchain Trilemma.}
\label{tab:trilemma-examples}
\end{table}

The trilemma is not a proven impossibility theorem---it is an empirical observation about the difficulty of simultaneously achieving all three goals. Much of current blockchain research (sharding, rollups, new consensus mechanisms) aims to push the boundaries of this trade-off, and a ``solution'' to the trilemma remains one of the holy grails of blockchain engineering.
